\section{Herkunft: wie $2/9$ überhaupt sichtbar wurde}
Bevor gerechnet wird, gehört die Frage geklärt, woher die zentrale Zahl
dieses Dokuments stammt -- denn sie wurde nicht gesucht, sondern gefunden.
Der Zugang war die Hilbertraum-Übersetzung von FFGFT. Die Brücke der
Dok.~230/231/232 bildet die Loop-Observablen des Korpus als
selbstadjungierte Operatoren auf einem faserwertigen Hilbertraum ab; erst als
diese Übersetzung bis zu den Massen fortgeführt wurde (Dok.~282), zeigte sich
die entscheidende Struktur: die drei Lepton-Massen sind das Spektrum eines
hermiteschen Z$_3$-Zirkulanten auf der $\mathbb{C}^3$-Faser des vollen Raums
$\mathcal{H}=L^2(T^4)\otimes\mathbb{C}^3$. In den Eigenwerten dieses
Operators erscheint der Phasenwinkel $\theta=2/9$ -- nicht als Ansatz, den
man einsetzt, sondern als Wert, den die Diagonalisierung ausgibt. $2/9$ ist
also \emph{kein} Startpunkt der Massenrechnung, sondern eine nachträglich
sichtbar gewordene Ordnung: erst der Wechsel der Darstellung -- von der
$\xi$-Leiter zur Operator-/Spektralsicht -- hat sie freigelegt. Diese
Gegenrechnung prüft daher nicht eine Eingabe, sondern eine \emph{Entdeckung};
das ist der Grund, warum die Richtung von den Daten her (unten) überhaupt
aussagekräftig ist. Der Zusammenhang zur $\xi$-Leiter, die die Massen ohne
$2/9$ berechnet, wird in der Schlusssektion \enquote{Zwei Schichten} geklärt.

\section{Zweck und Richtung der Rechnung}
Der FFGFT-Korpus rechnet grundsätzlich von $\xi=4/30000$ \emph{weg}: aus der
Geometrie folgen Relationen, die anschließend an Daten geprüft werden. Dieses
Dokument dreht die Richtung um. Ausgangspunkt sind ausschließlich die
empirischen Leptonmassen (PDG), und gefragt wird, was die FFGFT-Relationen
des Leptonsektors \emph{aus den Daten} machen -- welche Strukturzahlen die
Empirie zurückliefert, mit welchen Unsicherheiten, und wo Restspannungen
sitzen. Der Modus ist der Kompatibilitätsmodus: keine neue Herleitung, sondern
eine unabhängige Gegenrechnung nach P35, mit der Ebenentrennung nach P40
(exakt sind nur Verhältnisse; Absolutwerte tragen die Brückenkonstanten).
Alle Zahlen dieses Dokuments stammen aus
\texttt{python/Dok292\_Skripte/leptonen\_empirie\_check.py} (numpy-only, Seed
20780458) und sind mit einem Aufruf reproduzierbar.

Die Eingaben sind
\begin{align}
  m_e    &= 0{,}51099895069(16)\ \text{MeV},\\
  m_\mu  &= 105{,}6583755(23)\ \text{MeV},\\
  m_\tau &= 1776{,}86(12)\ \text{MeV}.
\end{align}
Die $\tau$-Masse ist um Größenordnungen die unsicherste der drei; sie
dominiert, wie sich zeigt, sämtliche Restspannungen.

\section{Ebene 1: der Zirkulant -- Koide-Relation und Winkel}

\subsection{Koide-Verhältnis aus den Daten}
Das Koide-Verhältnis $Q=(m_e+m_\mu+m_\tau)/(\sqrt{m_e}+\sqrt{m_\mu}+\sqrt{m_\tau})^2$
ergibt aus den PDG-Werten
\begin{equation}
  Q = 0{,}666660511 \pm 0{,}0000068,
\end{equation}
wobei die Unsicherheit fast vollständig von $m_\tau$ getragen wird. Die
Abweichung vom Exaktwert $2/3$ beträgt $-6{,}2\cdot10^{-6}$, das sind
$-0{,}91\,\sigma$. Die Daten sind mit exakt $2/3$ verträglich; sie
\emph{erzwingen} den Exaktwert nicht, aber sie widersprechen ihm auf
$5$--$6$ signifikanten Stellen nicht.

\subsection{Zirkulant-Fit: Amplitude und Winkel aus drei Massen}
Nach Dok.~282 sind die $\sqrt{m_k}$ Eigenwerte eines hermiteschen
Z$_3$-Zirkulanten, $\sqrt{m_k}=M\,(1+r\cos(\theta+2\pi k/3))$, mit den beiden
reinen Strukturzahlen $r=\sqrt2$ und $\theta=2/9$ und der Identität
$Q=(1+r^2/2)/3$. Die Fourier-Extraktion aus den drei Datenpunkten liefert
(Zuordnung $k=0,1,2\to\tau,e,\mu$; alle zyklischen Zuordnungen ergeben
denselben Winkel)
\begin{equation}
  M = 313{,}84\ \text{MeV},\qquad
  r_{\text{emp}} = 1{,}4142005,\qquad
  \theta_{\text{emp}} = 0{,}22222963 \pm 8{,}4\cdot10^{-6}.
\end{equation}
Die Amplitude weicht von $\sqrt2=1{,}4142136$ um $1{,}3\cdot10^{-5}$ ab --
über $Q=(1+r^2/2)/3$ ist das \emph{dieselbe} Aussage wie die
$Q$-Abweichung, keine unabhängige. Der Winkel liegt $+0{,}89\,\sigma$ über
$2/9=0{,}22222222$, ebenfalls $\tau$-limitiert.

\subsection{Der $\tau$-freie Test: Winkel aus $m_\mu/m_e$ allein}
Dok.~282 benennt $m_\mu/m_e$ als schärfsten Test, ob $\theta$ exakt $2/9$
ist: hält man $r=\sqrt2$ fest, kürzt sich $M$ im Verhältnis, und die beiden
präzisesten Massen der Physik bestimmen $\theta$ allein. Die Auflösung der
Gleichung
$\bigl(1+\sqrt2\cos(\theta+4\pi/3)\bigr)^2/\bigl(1+\sqrt2\cos(\theta+2\pi/3)\bigr)^2
= m_\mu/m_e$ ergibt
\begin{equation}
  \theta(\mu/e) = 0{,}2222220471,\qquad |\theta-2/9| = 1{,}75\cdot10^{-7}.
\end{equation}
Das reproduziert die Angabe aus Dok.~282 ($\approx1{,}8\cdot10^{-7}$, sieben
signifikante Stellen) unabhängig. Zwei Lesarten sind zu trennen. Gemessen an
der winzigen PDG-Unsicherheit von $m_\mu/m_e$ ist die Restdifferenz formal
hochsignifikant -- \emph{unter der Zusatzannahme}, dass $r$ exakt $\sqrt2$
ist. Lässt man $r$ über $Q$ mitlaufen, ist die Restdifferenz keine
unabhängige Spannung, sondern erneut dieselbe eine $\tau$-Spannung von
$\sim0{,}9\,\sigma$. Es gibt im Leptonsektor also nicht vier kleine
Abweichungen, sondern \emph{eine einzige}, die in vier Gestalten erscheint.

\subsection{Vorwärtsprobe und die eine Vorhersage}
Setzt man umgekehrt $r=\sqrt2$ und $\theta=2/9$ exakt, reproduziert der
Zirkulant $m_\mu/m_e$ auf $+0{,}0010\,\%$ und $m_\tau/m_\mu$ auf
$+0{,}0060\,\%$. Die Invertierung nach $m_\tau$ macht daraus eine
falsifizierbare Vorwärts-Aussage: aus $(m_e, m_\mu)$ und den beiden
Strukturzahlen folgt
\begin{equation}
  m_\tau^{\text{pred}} = 1776{,}968\ \text{MeV}
\end{equation}
(die drei Invertierungswege -- $Q=2/3$ exakt, $\theta=2/9$ exakt im
Drei-Massen-Fit, direkte Vorhersage aus $(m_e,m_\mu;\sqrt2,2/9)$ -- liefern
$1776{,}969$, $1776{,}967$ und $1776{,}968$ MeV, also denselben Wert). Gegen
PDG $1776{,}86\pm0{,}12$ MeV sind das $+0{,}90\,\sigma$. Eine
$\tau$-Massenmessung auf dem angestrebten Belle-II-Niveau von
$\sim0{,}05$ MeV entscheidet über diese Spannung: landet der Wert bei
$\approx1776{,}97$, ist die $\tau$-Richtung aufgelöst und das Paar
$(\sqrt2, 2/9)$ auf dem $10^{-5}$-Niveau bestätigt; landet er signifikant
darunter, ist die Näherung auf diesem Niveau widerlegt. Dies ist eine echte
Vorhersage im Sinne der Korpus-Methodologie, keine Retrodiktion: der Wert
stand vor der Messung fest und ist hier vorab gebucht. Was der Test
\emph{nicht} entscheidet, zeigt die Präzisionsanalyse in
Abschnitt~\ref{sec:praezision}.

\section{Ebene 2: die $\xi$-Leiter -- Verhältnisse und Reverse-$\xi$}

\subsection{Leiter-Verhältnisse gegen die Daten}
Die Leiterdarstellung $m_i = r_i\,\xi^{p_i}\,v$ (Dok.~006/046; Leptonen:
$r_e=4/3$, $p_e=3/2$; $r_\mu=16/5$, $p_\mu=1$; $r_\tau=25/9$, $p_\tau=2/3$
nach K2) macht die Verhältnisse parameterfrei:
\begin{align}
  \frac{m_\mu}{m_e}\Big|_{\text{Leiter}} &= \tfrac{12}{5}\,\xi^{-1/2} = 207{,}85
  \quad(\text{PDG } 206{,}77;\ +0{,}52\,\%),\\
  \frac{m_\tau}{m_\mu}\Big|_{\text{Leiter}} &= \tfrac{125}{144}\,\xi^{-1/3} = 16{,}992
  \quad(\text{PDG } 16{,}817;\ +1{,}04\,\%),\\
  \frac{m_\tau}{m_e}\Big|_{\text{Leiter}} &= \tfrac{25}{12}\,\xi^{-5/6} = 3532
  \quad(\text{PDG } 3477;\ +1{,}57\,\%).
\end{align}
Die Leiter trägt also die Größenordnungen und die $\sim1\,\%$-Genauigkeit --
nicht die Präzision des Zirkulanten. Beide Ebenen sind getrennt zu lesen, wie
in Dok.~282 und P40 deklariert: der Zirkulant liefert vier bis fünf Stellen
aus $(\sqrt2, 2/9)$, die Leiter liefert die Hierarchie aus $\xi$.

\subsection{Reverse-$\xi$: was die Daten für $\xi$ zurückgeben}
Löst man die drei Leiter-Verhältnisse nach $\xi$ auf, geben die Daten zurück
\begin{equation}
  \xi_{\mu/e} = 1{,}347\cdot10^{-4},\qquad
  \xi_{\tau/\mu} = 1{,}375\cdot10^{-4},\qquad
  \xi_{\tau/e} = 1{,}358\cdot10^{-4},
\end{equation}
also $+1{,}1$ bis $+3{,}2\,\%$ über $4/30000=1{,}3333\cdot10^{-4}$. Die drei
Werte streuen untereinander um $\sim2\,\%$ -- die Leiter ist in den
Verhältnissen auf dem Prozentniveau konsistent, nicht exakt. Das ist der
ehrliche empirische Status der Leiterdarstellung, unverändert gegenüber der
Deklaration in Dok.~006 (mittlere Abweichung unter $1\,\%$ auf
Massenebene; auf Verhältnisebene entsprechend $1$--$3\,\%$).

\subsection{Brückenkonstanten nach P40}
Pro Lepton ergibt die Auflösung nach der Brückenkonstante
$v_{\text{eff}}=m_i/(r_i\,\xi^{p_i})$:
\begin{equation}
  v_{\text{eff}}^{(e)} = 248{,}9\ \text{GeV},\qquad
  v_{\text{eff}}^{(\mu)} = 247{,}6\ \text{GeV},\qquad
  v_{\text{eff}}^{(\tau)} = 245{,}1\ \text{GeV},
\end{equation}
gestreut um den Higgs-Vakuumerwartungswert $246{,}22$ GeV
($+1{,}10/+0{,}58/-0{,}46\,\%$). Äquivalent: die erforderlichen
$r_{\text{eff}}$ bei festem $v$ weichen um dieselben Prozente von den
deklarierten rationalen Werten ab ($r_e$: $1{,}348$ statt $4/3$ usw.). Genau
das ist der P40-Befund von der Datenseite: die rationalen Strukturzahlen der
Leiter gelten auf dem Prozentniveau, die Restabweichung ist in die
Brückenkonstante absorbiert und in Verhältnissen \emph{nicht} vollständig
weggekürzt.

\subsection{$K_{\text{frak}}$-Gegenprobe: kein einfaches Muster}
Die naheliegende Frage, ob die Leiter-Restabweichungen dem fraktalen
Korrekturfaktor folgen, verneint die Rechnung: $K_{\text{frak}}=1-100\xi$
entspricht $-1{,}333\,\%$, die Halb- und Dreihalb-Potenzen $-0{,}665$ und
$-1{,}995\,\%$ -- die gemessenen Abweichungen $+0{,}52$ und $+1{,}04\,\%$
passen weder im Betrag noch im Vorzeichen auf eine einzelne
$K_{\text{frak}}$-Potenz. Auffällig ist allein, dass die
$\tau/\mu$-Abweichung ziemlich genau das Doppelte der $\mu/e$-Abweichung ist;
ob dahinter Struktur oder Zufall steht, bleibt nach P35 ungebucht. Der Punkt
wird als offen deklariert, nicht wegerklärt.

\subsection{Leiterübergreifende Kopplung: ein generationslineares Korrekturgesetz}
Bevor dieser Befund gebucht wird, ist die Buchungsebene festzulegen, denn
sie entscheidet über seine Zulässigkeit. Es gibt zwei einander
\emph{ausschließende} Rechenmodi. \emph{Modus~1} rechnet alle
Leptonverhältnisse rein aus $\xi$, ohne Referenzpunkt: dann sind $\mu/e$,
$\tau/\mu$ und $\tau/e$ allesamt echte Vorhersagen (jeweils $0{,}5$--$1{,}6\,\%$
neben den Daten), und alle drei sind legitime Prüfpunkte. \emph{Modus~2}
wendet P42 an: $m_\mu/m_e$ ist dann der deklarierte Anker, \emph{verbraucht}
und keine Prüfgröße mehr -- man darf seine Leiter-Abweichung nicht mehr
buchen, sonst nutzt man dieselbe Größe zweimal (einmal als Anker, einmal als
Test). Die folgende Kopplungsanalyse gehört \emph{ausschließlich} in
Modus~1, den referenzfreien; in Modus~2 ist sie gegenstandslos, weil dort
$\mu/e$ aus dem Vergleich fällt und als einzige unabhängige Leiter-Prüfung
$m_\tau$ bleibt -- die ohnehin über den Zirkulanten läuft, nicht über die
grobe Leiter. Innerhalb von Modus~1 gilt: die Aussage \enquote{kein
einfaches Muster} betrifft die Frage nach einer \emph{einzelnen}
$K_{\text{frak}}$-Potenz. Fragt man stattdessen, ob die
Leiterstufen \emph{untereinander} gekoppelt sind, zeigt sich sehr wohl
Struktur. Schreibt man die erforderliche Korrektur je Stelle als effektiven
Faktor $N_g$ über $K=1-N_g\,\xi$ (bzw.\ als Exponenten $p$ über
$K_{\text{frak}}^{p}$), so ergeben die Daten
\begin{align}
  \text{$\mu/e$:}\quad & N_1 = 38{,}89,\quad p_1 = 0{,}3873,\\
  \text{$\tau/\mu$:}\quad & N_2 = 77{,}06,\quad p_2 = 0{,}7694,\\
  \text{$\tau/e$:}\quad & N_3 = 115{,}55,\quad p_3 = 1{,}1567.
\end{align}
Zwei Beziehungen sind hier nicht zufällig. Erstens ist die
Exponenten-Struktur \emph{multiplikativ konsistent}: $p_3 = p_1+p_2$ auf
$6\cdot10^{-14}$ genau. Die $\tau/e$-Stelle trägt also \emph{keine} neue
Information -- sie ist das Produkt der beiden Stufen, wie es eine
durchgängige multiplikative Korrektur $K^{p_1}\cdot K^{p_2}=K^{p_1+p_2}$
verlangt. Zweitens stehen die effektiven Faktoren im Verhältnis
$N_1:N_2:N_3 = 1:1{,}981:2{,}971 \approx 1:2:3$: die Korrektur wächst
\emph{linear mit der Generationszahl}, $N_g \approx g\cdot N_0$ mit
$N_0 \approx 38{,}6$ (die drei Einzelschätzungen $N_1$, $N_2/2$, $N_3/3$
liegen bei $38{,}89$, $38{,}53$, $38{,}52$ -- untereinander auf unter
$1\,\%$ konsistent).

Das ist der strukturelle Grund, warum ein konstantes $K_{\text{frak}}$ die
Leiter nicht schließt und warum die Frage nach der \enquote{richtigen}
Zahl statt $100$ (etwa $89$) ins Leere läuft: die Leiter verlangt keinen
festen Faktor, sondern einen \emph{generationszählenden}, $N_g=g\cdot N_0$.
Bei $\alpha$ ist nur eine Skala im Spiel (das $\mu e$-Mittel) -- ein Faktor
genügt, und $\alpha$ passt. In der Leiter überlagern sich drei Generationen,
sodass $N$ linear wächst; ein konstanter Faktor kann das prinzipiell nicht
leisten. Damit bekommt die zuvor als Bringschuld benannte
\emph{generationsabhängige} Verkleidung eine konkrete Form: sie ist in
erster Näherung linear in $g$.

Die Toleranz-Bilanz zeigt, wie weit das Gesetz trägt (referenzfreier
Modus~1, gemeinsames $N_0=38{,}65$):
\begin{center}
\begin{tabular}{lrrrr}
\hline
Stelle & roh & roh [$\sigma_{\text{exp}}$] & korrigiert & korr.\ [$\sigma_{\text{exp}}$]\\
\hline
$\mu/e$    & $+0{,}521\,\%$ & $2{,}4\cdot10^{5}$ & $+0{,}0033\,\%$ & $1{,}5\cdot10^{3}$\\
$\tau/\mu$ & $+1{,}038\,\%$ & $154$              & $-0{,}0031\,\%$ & $0{,}47$\\
$\tau/e$   & $+1{,}565\,\%$ & $232$              & $-0{,}0052\,\%$ & $0{,}78$\\
\hline
\end{tabular}
\end{center}
Das Gesetz drückt die Leiter-Abweichung um rund den Faktor $300$, von
$0{,}5$--$1{,}6\,\%$ auf $\sim5\cdot10^{-3}\,\%$. Die beiden
$\tau$-Stellen fallen damit \emph{in} ihre Messtoleranz ($<1\,\sigma$) --
für sie ist das Gesetz praktisch vollständig. Der $\mu/e$-Eintrag bleibt bei
$\sim1{,}5\cdot10^{3}\,\sigma$, nicht weil die Korrektur dort versagt,
sondern weil $\mu/e$ rund $500$-mal präziser gemessen ist: schon ein
$0{,}003\,\%$-Rest ist in $\sigma$-Einheiten riesig. Genau hier greift
aber die Modus-Grenze: dieser $\mu/e$-Eintrag ist \emph{nur} im
referenzfreien Modus~1 überhaupt eine Aussage; unter P42 ist $\mu/e$ der
Anker und fällt aus der Bilanz. Der verbleibende Rest der beiden
$\tau$-Stellen sitzt in der Feinabweichung des Verhältnisses vom exakten
$1:2:3$ ($\sim1$--$2\,\%$) und in der noch unbestimmten Grundeinheit $N_0$
-- das ist die nächste offene Rechnung, nicht ein Widerspruch.

Offen -- und nach P35 ausdrücklich nicht wegerklärt -- bleibt zweierlei. Die
Grundeinheit $N_0\approx38{,}6$ ist noch nicht aus der Geometrie
hergeleitet; naheliegende Kandidaten ($100/e=36{,}8$, $100/\varphi^2=38{,}2$,
glatt $39$) treffen sie nicht genau genug, um behauptet zu werden. Und der
Phasensektor -- der $\delta$-Feinwert aus Dok.~291 mit $p_\delta=-0{,}88$ --
fügt sich \emph{nicht} offensichtlich in dieselbe generationslineare
Magnitude-Leiter; ob Magnitude- und Phasensektor über ein gemeinsames Gesetz
koppeln, ist eine eigene, hier nicht beantwortete Frage. Der Befund ist also
nicht die Lösung des Leiter-Restes, sondern seine Präzisierung: aus
\enquote{drei krumme Exponenten} wird \enquote{ein generationslineares
Gesetz $N_g=g\cdot N_0$ mit noch unbestimmtem $N_0$}.

Ausdrücklich festzuhalten ist die Grenze der Aussagekraft, die aus der
Modus-Trennung folgt. Das $1:2:3$-Muster ist an drei Punkten belegt
\emph{nur} im referenzfreien Modus~1. Sobald man P42 anwendet und $m_\mu/m_e$
als Anker deklariert (die für die Zirkulant-Präzision und die
$m_\tau$-Vorhersage nötige Ebene), fällt $\mu/e$ als Prüfpunkt weg, und das
Gesetz stützt sich nur noch auf $\tau/\mu$ und $\tau/e$ -- die dieselbe
$\tau$-Masse teilen und deren Präzision gering ist. Die scheinbar dreifache
Bestätigung ist damit an die Referenzfreiheit gebunden. Beides zugleich --
$\mu/e$ als Anker \emph{und} als $N_g$-Prüfpunkt -- ist unzulässige
Doppelbuchung. Der Zirkulant, der $\mu/e$ bereits auf $10^{-5}$ trägt (Weg~A
in Abschnitt~\ref{sec:praezision}), braucht die $g$-Korrektur ohnehin nicht;
sie ist eine Aussage allein über die grobe $\xi$-Leiter im referenzfreien
Modus, nicht über die präzise Zirkulant-Schicht.

Ein häufiges Missverständnis ist hier auszuräumen: die $100$ aus
$K_{\text{frak}}=1-100\,\xi$ und die Grundeinheit $N_0\approx38{,}6$ sind
\emph{nicht} dieselbe Größe, und die Leiter braucht darum keineswegs
\enquote{weniger Rekursionen} als der $\alpha$-Sektor. Die $100$ ist die
eine feste Skala einer korpusweit wirkenden Konstante (X-Flip,
$T_{\text{CMB}}$, $E_0$-Brücke); $N_0$ ist der Faktor \emph{pro Generation}
in einem anderen Sektor. Der Unterschied ist nicht \emph{weniger}, sondern
\emph{anders}: kein fester Faktor, sondern ein generationszählender
$N_g=g\cdot N_0$ -- bei $g=2$ bereits $77$, bei $g=3$ bereits $116$, also
teils \emph{mehr} als $100$. Ein Vergleich der beiden Zahlen als
\enquote{Rekursionstiefen} geht daher ins Leere.

Bei der Bewertung möglicher Ausdrücke für $N_0$ ist der richtige Maßstab
nicht die punktgenaue Übereinstimmung -- die ist prinzipiell unerreichbar --,
sondern die Verträglichkeit \emph{innerhalb der Toleranz}, mit der $N_0$
überhaupt bestimmbar ist. Diese Toleranz ist beträchtlich: $N_0$ folgt im
referenzierten Modus allein aus den $\tau$-Stellen, und deren Genauigkeit
ist durch die $\tau$-Masse ($\pm0{,}12$ MeV) begrenzt. Propagiert man das,
ergibt sich $N_0 = 38{,}52 \pm 0{,}21$, also eine relative Bestimmungs\-un\-si\-cher\-heit
von rund $0{,}5\,\%$. Ein Kandidatenausdruck ist demnach nicht daran zu
messen, ob er $38{,}6$ auf der Nachkommastelle trifft, sondern ob er in
dieses Intervall fällt.

Unter diesem Maßstab -- und weiterhin als \emph{unbewiesene} Beobachtung
nach P35, nicht als Herleitung -- ist $100/\varphi^2 = 38{,}20$ ein
\emph{zulässiger} Kandidat: der Abstand zu $N_0$ beträgt $1{,}6\,\sigma_{N_0}$,
liegt also innerhalb der Toleranz. Als Punktwert verfehlt der Ausdruck die
$38{,}6$; innerhalb der Toleranz, in der $N_0$ bekannt ist, trifft er. Wäre
$N_0=100/\varphi^2$, so folgte die Leiter-Grundeinheit über einen einzigen
Faktor $\varphi^{-2}$ aus der $K_{\text{frak}}$-Skala -- Magnitude- und
Korrektursektor hingen an derselben Konstante. Der Befund ist damit nicht
\enquote{knapp verfehlt und verworfen}, sondern \enquote{mit der
gegenwärtigen Bestimmung verträglich, nicht widerlegt}. Eine schärfere
$\tau$-Masse (Belle~II) oder eine geometrische Herleitung von $N_0$ wird
zwischen $100/\varphi^2$ und den Alternativen entscheiden; bis dahin bleibt
es ein zulässiger, nicht ein bewiesener Ausdruck.

\section{Relation zu den Messgenauigkeiten}
\label{sec:praezision}
Die bisherigen $\sigma$-Angaben beziehen sich je auf eine Größe. Die
vollständige Frage lautet: Wie liegen die FFGFT-Strukturzahlen relativ zur
\emph{gesamten} Präzisionsstruktur der Daten? Sie wird durch eine
Zwei-Parameter-Analyse beantwortet. Die zwei unabhängigen Verhältnisse
$R_1=m_\mu/m_e$ und $R_2=m_\tau/m_\mu$ bestimmen das Paar $(r,\theta)$
exakt; ihre relativen Messgenauigkeiten sind jedoch dramatisch verschieden:
$\sigma_{R_1}/R_1 = 2{,}2\cdot10^{-8}$ gegen
$\sigma_{R_2}/R_2 = 6{,}8\cdot10^{-5}$ -- das Verhältnis $\mu/e$ ist rund
$3100$-mal präziser gemessen. Die Fehlerfortpflanzung auf $(r,\theta)$
ergibt daher eine extrem anisotrope Kovarianz: eine \emph{dünne}
Hauptrichtung mit $\sigma \approx 3{,}4\cdot10^{-10}$ (von $\mu/e$
getragen) und eine \emph{dicke} mit $\sigma \approx 1{,}7\cdot10^{-5}$
(von $\tau$ getragen); die Korrelation zwischen $r$ und $\theta$ beträgt
praktisch $-1$.

In diesem Bild liegt der Punkt $(\sqrt2, 2/9)$ vom Daten-Fit
$(1{,}41420051,\ 0{,}22222963)$ entfernt um $0{,}90\,\sigma$ in der
dicken Richtung -- das ist die eine $\tau$-Spannung aus
Abschnitt~2 -- und um rund $450\,\sigma$ in der dünnen. Ausformuliert:
der Zirkulant mit \emph{exakt} $(\sqrt2, 2/9)$ verfehlt $m_\mu/m_e$ um
$1{,}0\cdot10^{-5}$ relativ, während die Messung auf $2{,}2\cdot10^{-8}$
steht. Als exaktes Paar sind die beiden Strukturzahlen von den heutigen
Daten also bereits ausgeschlossen; bestätigt sind sie auf dem
$10^{-5}$-Niveau -- genau die \enquote{vier bis fünf Stellen}, die
Dok.~282 deklariert, hier als Messrelation präzisiert. Daraus folgt auch
eine Grenze der $m_\tau$-Vorhersage: ein perfektes Belle-II-Ergebnis löst
nur die dicke $\tau$-Richtung; der Rest von $1{,}75\cdot10^{-7}$ in
$\theta$ (äquivalent $10^{-5}$ in $R_1$) ist ein \emph{heute schon
gemessener}, stehender Offset, den jede Exaktheitsbehauptung erklären
müsste -- naheliegende Kandidaten wären radiative Korrekturen an den
Polmassen, wobei deren typische Größe ($\alpha/\pi \sim 10^{-3}$) das
Rätsel eher vergrößert, warum die Übereinstimmung überhaupt bis
$10^{-5}$ reicht. Der Punkt bleibt nach P35 offen gebucht.

\subsection{Gedankenmodell: die Umkehrung der Präzisionsaussage}
Der $450\,\sigma$-Befund ruht vollständig auf dem Vertrauen in die
angegebene Präzision von $m_\mu/m_e$. Die Umkehrung -- die Präzision
könnte nicht das bedeuten, was sie behauptet -- ist als Gedankenmodell
zulässig und kommt in drei Stärkegraden. \emph{Schwach:} die
experimentelle Unsicherheit wäre um den Faktor $\sim450$ unterschätzt --
quantitativ kaum haltbar, da dieselbe Präzision in einem Netz unabhängiger
Kreuzchecks (Myonium-Spektroskopie, $g\!-\!2$, QED-Tests) hängt.
\emph{Mittel:} $m_\mu/m_e$ ist keine abgelesene, sondern eine durch
gebundene QED \emph{extrahierte} Größe; die Präzisionsangabe setzt die
Vollständigkeit der Extraktionstheorie voraus, und ein fehlender Term bei
$10^{-5}$ würde den Konflikt auflösen -- $450$-mal über dem geschätzten
Theoriefehler, grenzwertig, aber logisch offen. \emph{Stark (und
FFGFT-nativ):} die Präzision ist echt, gilt aber für die \emph{angezogene}
Größe -- die Polmassen im SI-Rahmen --, während die Zirkulant-Relation auf
der geometrischen Ebene darunter läge. Dann widerspricht nichts einander:
der $10^{-5}$-Offset wäre kein Fehler, sondern die \emph{gemessene Dicke
der Verkleidungsschicht} im $\mu/e$-Kanal. Das ist die konsequente
Verlängerung von P40: in Verhältnissen kürzt sich nur der
generations\emph{unabhängige} Teil der fraktalen/SI-Korrektur; ein
generationsabhängiger Rest bleibt stehen -- und hätte genau die Gestalt
dieses Offsets. Der Preis der starken Lesart ist eine Bringschuld: der
Offset müsste aus der Verkleidung berechenbar sein (generationsabhängige
radiative Struktur), sonst ist die Lesart unfalsifizierbar. Sie wird hier
als Gedankenmodell gebucht, nicht als Auflösung; der Widerspruchs-Befund
des vorigen Abschnitts bleibt für die Polmassen-Ebene bestehen.

\subsection{Aufgelöst durch Deklaration: der Referenzpunkt $m_\mu/m_e$ (P42)}
Der Rahmen entscheidet die Frage nicht über eines der Gedankenmodelle,
sondern über die Anspruchsschicht: FFGFT bleibt grundsätzlich bei der rein
theoretischen Ableitung aus $\xi$, und der Leptonsektor erhält -- analog zur
kosmischen Skala (P39) -- genau \emph{einen} deklarierten Referenzpunkt,
das Verhältnis $m_\mu/m_e$. Bei strukturellem $r=\sqrt2$ (aus $Q=2/3$)
fixiert es den operativen Winkel $\theta_{\text{ref}}=0{,}2222220471$;
die Strukturzahl $2/9$ ist dessen rationale Adresse auf sieben Stellen,
kein Exaktheitsanspruch auf Polmassen-Ebene. Damit ist der
$450\,\sigma$-Befund keine Falsifikation mehr, sondern Teil der
Kalibrierung. Zwei Abgrenzungen: der Status von $\xi$ ist davon unberührt
(FFGFT behauptet $\xi$ als zwingend aus der Geometrie; dieser Streitpunkt
ist hier nicht verhandelt), und die Leptonmassen sind an sich keine Eingabe
der Theorie -- der Referenzpunkt gehört zur Vergleichsebene, er verortet
die geometrische Relation in den Polmassen-Daten und wäre ohne
Vergleichsabsicht nicht nötig. Der ehrliche Preis liegt in der Testbilanz:
$m_\mu/m_e$ wird als Anker verbraucht und scheidet als Prüfgröße aus;
Prüfgröße ist allein $m_\tau$. Der Gewinn ist eine eingabefehlerfrei
scharfe Vorhersage:
\begin{equation}
  m_\tau = 1776{,}9690 \pm 0{,}0001\ \text{MeV},
\end{equation}
wobei die Unsicherheit allein aus $(m_e, m_\mu)$ stammt und praktisch
verschwindet; gegen PDG bleibt die eine $+0{,}91\,\sigma$-Spannung, und
das Belle-II-Weltmittel entscheidet ohne jede Ausweichmöglichkeit. Für die
Brückenfrage (Dok.~291) verschärft die Deklaration das Ziel: die Dynamik
muss $\theta_{\text{ref}}$ selektieren; $2/9$ bleibt dessen Adresse und
das vorregistrierte Stage-10-Ziel, dessen Orbit-Auflösung ($\sim0{,}1$)
vom $10^{-7}$-Unterschied unberührt ist. Gebucht als P42 in Dok.~190.

Dieselbe Betrachtung stellt die $\xi$-Leiter ein: ihre Restabweichungen
entsprechen $2{,}4\cdot10^{5}$ ($\mu/e$), $154$ ($\tau/\mu$) und $232$
($\tau/e$) experimentellen Standardabweichungen. Die Leiter-Residuen sind
also nicht datenlimitiert, sondern theorieseitig (P40-Brückenkonstanten);
die Messpräzision hat die Leiterdarstellung um zwei bis fünf
Größenordnungen überholt. Für die $m_\tau$-Vorhersage schließlich beträgt
die Testschärfe (Trennung Vorhersage gegen heutigen Zentralwert) bei
$\sigma_{m_\tau}=0{,}12/0{,}05/0{,}02$ MeV entsprechend
$0{,}9/2{,}2/5{,}4\,\sigma$ -- eine einzelne Belle-II-Messung ist
indikativ, erst das kombinierte Weltmittel entscheidend.

\section{Die Feinstrukturkonstante aus den Leptonmassen}
Die Feinstrukturkonstante gehört in diese Gegenrechnung, weil sie in FFGFT
eine Leptonmassen-Aussage ist: Dok.~011 setzt
$\alpha = \xi\,(E_0/\text{MeV})^2$ mit der strukturellen Energieskala
$E_0=\sqrt{m_e\,m_\mu}$ und der deklarierten fraktalen Korrektur,
$\alpha^{-1} = (7500/E_0^2)\,K_{\text{frak}}$. Von der Datenseite: das
gemessene $\alpha$ verlangt $E_0^{\text{req}}=7{,}39798$ MeV, das
geometrische Mittel liefert $7{,}34788$ MeV -- Verhältnis $1{,}006819$,
gegen $K_{\text{frak}}^{-1/2}=1{,}006734$. Der deklarierte
Korrekturfaktor schließt die Lücke also bis auf $8\cdot10^{-5}$; die
geschlossene, gänzlich parameterfreie Relation
\begin{equation}
  \alpha \;=\; \frac{\xi\, m_e\, m_\mu}{K_{\text{frak}}\ \text{MeV}^2},
  \qquad K_{\text{frak}}=1-100\,\xi,
\end{equation}
liefert $1/\alpha = 137{,}059$ gegen gemessen $137{,}035999$.

Entscheidend ist hier, dass $\alpha$ \emph{nicht} über eine gefittete
Brückenkonstante läuft, sondern über \emph{zwei unabhängige Wege} zur
charakteristischen Energie $E_0$, die einander bestätigen (Dok.~011). Weg~1
ist empirisch: $E_0=\sqrt{m_e m_\mu}=7{,}348$ MeV, das geometrische Mittel
der gemessenen Massen. Weg~2 ist eigenständig-geometrisch:
$E_0^2=4\sqrt2\,m_\mu/\xi^{p}$ konstruiert $E_0=7{,}398$ MeV aus $\xi$
und $m_\mu$ allein, ohne das geometrische Mittel zu verwenden -- keine
Rückrechnung aus $\alpha$. Zwei Konsistenzen, die nichts voneinander
wissen, treffen zusammen: (i) Weg~2 trifft $7{,}398$ und liefert damit
$\alpha=\xi\,(E_0/\text{MeV})^2$ auf $+5\cdot10^{-6}$; (ii) die Differenz
der beiden Wege, $7{,}398/7{,}348=1{,}00682$, ist auf $9\cdot10^{-5}$ genau
$K_{\text{frak}}^{-1/2}$. Damit ist die $\alpha$-Stelle \emph{keine}
Kalibrierung, sondern überbestimmt: zwei getrennte Herleitungen treffen sich
auf $8\cdot10^{-5}$, und die Rekursionstiefe der Korrektur ist nicht frei,
sondern strukturell fixiert -- $E_0$ geht quadratisch in $\alpha$ ein,
weshalb die $K_{\text{frak}}^{1/2}$-Brücke auf $E_0$-Ebene einer vollen
$K_{\text{frak}}$-Ordnung auf $\alpha$-Ebene entspricht. Der \emph{gemessene}
Restwert $\alpha$ selbst wird von Weg~2 auf $5\cdot10^{-6}$ getroffen; dass
\emph{derselbe} Faktor $K_{\text{frak}}=1-100\,\xi$, der im Korpus den
gedämpften X-Flip und die $T_{\text{CMB}}$-Korrektur trägt, genau die Lücke
der beiden $E_0$-Wege schließt, ist eine sektorübergreifende Konsistenz --
kein Fit.

\subsection{Der $\alpha$-Restwert als zweiter Zeuge der Verkleidungsschicht}
Die Überbestimmung von $\alpha$ erlaubt eine Aussage, die der $\theta$-Fall
(Abschnitt~\ref{sec:praezision}) nur andeutet. Nimmt man beide Anker als
exakt -- das auf $1{,}5\cdot10^{-10}$ gemessene $\alpha$ und die
rein geometrische Weg-2-Konstruktion $E_0^2=4\sqrt2\,m_\mu/\xi^{p}$, die
keinen Messfehler trägt --, dann ist in $\alpha=\xi\,E_0^2$ die einzige
Größe mit Spielraum der empirische Weg~1, also $E_0=\sqrt{m_e m_\mu}$ und
damit die Polmassen selbst. Die $K_{\text{frak}}$-Brücke schließt $99{,}99\,\%$
der Lücke zwischen Weg~1 und dem von $\alpha$ verlangten $E_0$; es bleibt
ein Rest von $8{,}4\cdot10^{-5}$. Gemessen an der Präzision der Massen
($1{,}1\cdot10^{-8}$ auf $E_0$-Ebene) ist dieser Rest $\sim7700\,\sigma$
groß -- er ist zu groß, um Massen-Messrauschen zu sein, und damit eine
Aussage \emph{über} die Polmassen: unter der starken Lesart (Weg~2 exakt)
ist $\sqrt{m_e m_\mu}$ nicht die geometrisch exakt passende Größe.

Entscheidend ist, dass dieser $\alpha$-Rest \emph{modusunabhängig} gilt --
und zwar genau wegen der zwei Wege. Weg~2 benutzt $m_e$ \emph{nicht}: er
konstruiert $E_0$ aus $\xi$ und $m_\mu$ allein. Damit ist $m_\mu/m_e$ hier
\emph{kein} Anker; $m_e$ wird nirgends hineingesteckt und wieder
herausgeholt. Die Übereinstimmung von Weg~2 (ohne $m_e$) mit Weg~1 (mit
$m_e$) ist deshalb eine echte Vorhersage über $m_e$, und der Rest von
$8\cdot10^{-5}$ ist keine Selbstreferenz, sondern ein echter Test der
Polmasse gegen die Geometrie -- er bleibt eine Aussage, gleich ob man
referenzfrei oder mit P42 rechnet.

Damit lässt sich der Befund präzise einordnen -- schärfer, als eine
symmetrische \enquote{zwei Zeugen}-Formulierung es täte. Der Vergleich mit
der $\theta$-Seite (Abschnitt~\ref{sec:praezision}) zeigt eine
\emph{Asymmetrie}: dort wird $\theta_{\text{ref}}$ aus $m_\mu/m_e$
gewonnen, also ist $\mu/e$ Anker, und die $\sim450\,\sigma$ im
$\mu/e$-Kanal sind zur Hälfte Selbstreferenz -- ein Zeuge \emph{nur} im
referenzfreien Modus. Der $\alpha$-Rest hingegen ist \emph{modusunabhängig},
weil Weg~2 den zweiten Anker ($m_e$) gar nicht verwendet. Der starke,
in jedem Modus gültige Beleg für eine Verkleidungsschicht zwischen
geometrisch fundamentaler Größe und QED-angezogener Polmasse ist also
$\alpha$ allein; die $\theta$-Seite stützt denselben Schluss, aber nur unter
Referenzfreiheit. Der ehrliche Preis bleibt: die Aussage setzt die starke
Lesart voraus (Weg~2 als exakt), und die Bringschuld -- die Verkleidung zu
\emph{berechnen} statt sie als Rest zu benennen -- ist offen. Als
Kalibrierungsartefakt lässt sich der $\alpha$-Rest jedoch nicht abtun: dafür
ist $\alpha$ überbestimmt und die Massenpräzision zu hoch.

Konzeptionell -- und rechnerisch fast folgenlos, aber für die Architektur
zentral -- gilt: im erweiterten natürlichen Einheitensystem ist
$\alpha=1$ (Dok.~011/014, Ladung dimensionslos); der SI-Wert
$1/137{,}036$ ist Buchhaltung der Einheitenübersetzung, keine
fundamentale Kopplung. Dieselbe Architektur trägt die Gravitation:
$G=\xi^2/(4 m_{\text{char}})$ mal Umrechnung (Dok.~012/013) macht $G$
zur Buchhaltungsgröße zur SI -- die Reproduktion des SI-Werts ist per P40
Kalibrierung, testbar allein die Vorwärtskette. Der strukturelle Gewinn
liegt woanders: eine Kopplung, die nur Buchhaltung ist, muss nicht
quantisiert werden -- das Quantengravitationsproblem, an dem die meisten
Erweiterungen des Standardmodells scheitern, stellt sich in dieser
Architektur nicht.

\section{Zwei Schichten, die einander nicht brauchen: was $2/9$ leistet und was die SI-Umrechnung leistet}
Zwei Fragen, die leicht vermischt werden, sind hier zu trennen.

\emph{Erstens: Wird $2/9$ zur Massenberechnung benötigt?} Nein. Der
historische und weiterhin gültige Rechenweg ist die $\xi$-Leiter
(Dok.~006/046): $m_i=r_i\,\xi^{p_i}\,v$ mit $r\in\{4/3,16/5,25/9\}$ und
$p\in\{3/2,1,2/3\}$ liefert die Massen aus $\xi$ und den Leiter-Zahlen
allein, ohne dass $2/9$ irgendwo eingeht. Die Strukturzahl $2/9$ erscheint
erst ab Dok.~282 und in anderer Funktion: nicht als Eingabe, die die Massen
erzeugt, sondern als \emph{nachträglich entdeckte Ordnung} in ihnen -- die
drei $\sqrt{m_k}$ sind Eigenwerte eines Z$_3$-Zirkulanten mit Phasenwinkel
$2/9$. $2/9$ ersetzt den früher gefitteten Leiter-Koeffizienten durch eine
reine Zahl; es ist Verdichtung und Erklärung, keine Voraussetzung. Man
könnte $2/9$ aus dem Massensektor streichen, und die Leiter rechnete
weiter -- verloren ginge die \emph{Einsicht}, warum $Q$ exakt $2/3$ ist,
nicht die Zahlen. $2/9$ gehört also zur Erklärungs-, nicht zur
Rechenschicht.

\emph{Zweitens: Bestätigt diese Gegenrechnung, dass die SI-Umrechnung
notwendig und richtig ist?} In zwei getrennten Graden. Als
\emph{notwendig} bestätigt sie sich klar: jede Absolutzahl ($m_\tau$ in
MeV, $\alpha=1/137$, $G$) verlangt die Brückenkonstante, um in SI zu
landen, während die Verhältnisse ($m_\mu/m_e$) sie nicht brauchen. Genau
diese Zwei-Ebenen-Struktur ist P40, und der $450\,\sigma$-Befund ist ihr
schärfster Beleg: er zeigt, dass die gemessene Polmassen-Ebene (SI-definiert,
QED-angezogen) eine \emph{andere} Ebene ist als die geometrische -- ohne
den Umrechnungsschritt gäbe es die Differenz nicht. Als \emph{auf
Messpräzision validiert richtig} bestätigt sie sich hingegen \emph{nicht}:
die Umrechnung schließt bei $\alpha$ rund $98{,}8\,\%$ der Lücke (über
denselben Faktor $K_{\text{frak}}$, der auch in anderen Sektoren trägt),
aber der Rest von $\sim10^{-4}$ bleibt offen. Die SI-Umrechnung ist damit
strukturell notwendig und über Sektoren hinweg konsistent, aber sie ist
Kalibrierung, kein auf Nachkommastellen bestandener Test.

Der eigentliche Befund liegt in der Unabhängigkeit: die Leiter rechnet
Massen ohne $2/9$, der Zirkulant erklärt $Q$ ohne die Leiter, die
SI-Umrechnung übersetzt ohne $2/9$. Dass alle drei dennoch dasselbe $\xi$
und denselben Faktor $K_{\text{frak}}$ teilen, ist die Konsistenz des
Rahmens -- keine Abhängigkeit einer Schicht von der anderen.

\section{Einordnung}
Die Gegenrechnung von der Datenseite bestätigt die Zwei-Ebenen-Architektur
des Korpus, ohne ihr etwas hinzuzufügen oder zu ersparen. Auf der
Zirkulant-Ebene geben die Daten die beiden Strukturzahlen auf dem
$10^{-5}$-Niveau zurück; die Präzisionsanalyse trennt dabei zwei Befunde:
in der $\tau$-limitierten (dicken) Richtung eine einzige
$0{,}9\,\sigma$-Spannung, die die vorab gebuchte Vorhersage
$m_\tau = 1776{,}968$ MeV trägt -- und in der $\mu/e$-präzisen (dünnen)
Richtung einen heute schon gemessenen $10^{-5}$-Offset, der die
\emph{Exaktheit} des Paars $(\sqrt2, 2/9)$ bereits ausschließt und als
offene Feinstruktur gebucht ist. Auf der Leiter-Ebene
geben die Daten $\xi$ auf $1$--$3\,\%$ zurück, mit ehrlich offenen
Restabweichungen ohne erkennbares $K_{\text{frak}}$-Muster. Was diese
Rechnung \emph{nicht} zeigt: keine Herleitung von $\theta=2/9$ (der
dynamische Ort bleibt die offene Frage der Brückenarbeit), keine
Verschärfung der Leiter, keine neue Physik. Sie zeigt, dass die deklarierten
Ansprüche beider Ebenen von den PDG-Daten aus unabhängig reproduzierbar sind
-- in der Richtung, in der sich nichts anpassen lässt.

\subsection{Der Rest als Signatur der Unvollständigkeit}
Die Gegenrechnung hat die Aussagen geschärft -- von \enquote{zwei reine
Zahlen, kein Fit} zu vermessenen Größen mit bekannten Nachkommastellen und
bekannten Grenzen. Doch auf mehreren Ebenen bleibt ein Rest: der
$10^{-7}$-Offset in $\theta$, die $1$--$3\,\%$ der Leiter, der offene
$\delta$-Feinwert (Dok.~291), das über die Leiter-Stellen hinweg nicht
schließende $K_{\text{frak}}$-Muster. (Die $\alpha$-Stelle gehört
\emph{nicht} hierher: sie ist über zwei unabhängige $E_0$-Wege
überbestimmt und ihre Rekursionstiefe ist strukturell fixiert -- siehe
oben.) Keiner dieser Reste ist ein
\emph{Widerspruch} -- die Theorie stimmt überall auf ihrer jeweiligen
Genauigkeit --, aber keiner ist \emph{null}. Ein Residuum, das kleiner
wird, je genauer man hinsieht, ohne zu verschwinden, ist die Signatur einer
Näherung an etwas Tieferes, nicht die eines abgeschlossenen Fundaments.
Deshalb ist der Rest kein Fehler, sondern der ehrliche Beleg, dass die
Theorie \emph{nicht vollständig} ist.

Bemerkenswert ist, \emph{wo} er sitzt. Der geometrische Kern trägt keinen
Rest: $Q=2/3$ ist rein geometrisch und praktisch exakt; im erweiterten
Natursystem ist $\alpha=1$; die Leiter-Verhältnisse sind rationale Zahlen.
Der Rest entsteht durchgängig \emph{erst} beim Übergang ins Gemessene: der
$450\,\sigma$-Bruch erst auf der Polmassen-Ebene, die Leiter-Abweichung
erst in der Brückenkonstante. Die $\alpha$-Umrechnung ist die Ausnahme, die
die Regel schärft: dort fixieren \emph{zwei} unabhängige $E_0$-Wege die
Rekursionstiefe, sodass die Übersetzung nicht frei bleibt -- genau deshalb
ist $\alpha$ die eine Stelle ohne offenen Rest in diesem Sinne. Wo hingegen
nur \emph{ein} Weg vorliegt (die Leiter-Verhältnisse), bleibt die Tiefe der
Korrektur unbestimmt, und dort sitzt der Rest. Der Rest lebt also nicht im geometrischen Kern, sondern in
der \emph{Übersetzungsschicht} -- systematisch, nicht zufällig verteilt. Das
lokalisiert die Unvollständigkeit präzise: nicht in $\xi$, nicht im
Zirkulanten, sondern in der noch nicht \emph{hergeleiteten} fraktalen bzw.\
QED-Verkleidung. Vollständig wäre die Theorie erst, wenn sie diesen Rest
\emph{berechnet} statt ihn in eine Brückenkonstante zu absorbieren -- wenn
$K_{\text{frak}}$ nicht nur passt, sondern generationsabhängig aus der
Geometrie folgt. Der Maßstab dafür ist ausdrücklich \emph{nicht} die
punktgenaue Landung auf der letzten Nachkommastelle: eine geometrische
Herleitung soll nicht die QED-angezogene Polmasse Ziffer für Ziffer treffen
-- das wäre weder möglich noch das Ziel --, sondern die zugrunde liegende
Struktur so weit festlegen, dass die Vorhersagen \emph{innerhalb der
Toleranzen} bleiben, in denen die Größen überhaupt bestimmbar sind.
Vollständigkeit heißt hier: kein Rest mehr, der \emph{über} diese Toleranzen
hinausragt und nur durch einen freien Parameter geschlossen wird. Bis dahin
ist jeder solche Rest der Zeiger auf die nächste offene Herleitung, an der
die Vollständigkeit -- im Sinne der Toleranz, nicht der letzten Ziffer -- zu
messen sein wird.
